\chapter{Semântica operacional para uma linguagem imperativa simples}

\section{Introdução}

No capítulo anterior, definimos a semântica operacional da linguagem
\emph{Line}, que inclui atribuição, leitura e impressão. Com isso, já
conseguimos modelar programas que executam uma sequência fixa de
comandos. No entanto, programas reais precisam de \textbf{controle de
fluxo}: executar blocos diferentes dependendo de uma condição, e
repetir um bloco enquanto uma condição permanecer verdadeira.

Este capítulo estende a linguagem com \textbf{condicionais}
(if/else) e \textbf{repetição}
(while), formando a linguagem \emph{While}.
Para suportar condições, a linguagem de expressões também é estendida
com \textbf{valores booleanos} e operações de comparação. Definimos a
semântica small-step e big-step para todas as novas construções e
mostramos como essa especificação se traduz no interpretador
disponível em
\begin{flushleft}
  \verb|While/Interp/WhileInterp.hs|
\end{flushleft}

\section{A Linguagem While}

A linguagem \emph{While} estende \emph{Line} com dois novos comandos
e uma classe de expressões booleanas. A sua sintaxe abstrata é:

\[
  \begin{array}{lcl}
    \mathit{prog} & ::= & s^\star\\
    s & ::= & x \mathbin{:=} e
    \mid \mathbf{read}\; x
    \mid \mathbf{print}\; e \\
    &    & \mid\; \mathbf{if}\; e\; \mathbf{then}\; \{s^\star\}\;
    \mathbf{else}\; \{s^\star\} \\
    &    & \mid\; \mathbf{while}\; e\; \mathbf{do}\; \{s^\star\} \\
    e & ::= & v \mid x
    \mid e_1 + e_2 \mid e_1 \star e_2 \\
    &    & \mid\; e_1 < e_2 \mid e_1 = e_2 \mid e_1 \mathbin{|} e_2
    \mid \neg e \\
    v & ::= & n \mid \mathbf{true} \mid \mathbf{false}
  \end{array}
\]

Os valores agora incluem inteiros e booleanos. Em Haskell:

\begin{minted}{haskell}
data Value
  = VInt  Int
  | VBool Bool

data Exp
  = EValue Value
  | EVar   Var
  | Exp :+: Exp  -- adição
  | Exp :*: Exp  -- multiplicação
  | Exp :<: Exp  -- menor que
  | Exp :=: Exp  -- igualdade
  | Exp :|: Exp  -- disjunção
  | ENot Exp     -- negação

data Stmt
  = SAssign Var Exp
  | SRead   Var
  | SPrint  Exp
  | SIf    Exp Block Block
  | SWhile Exp Block

type Block = [Stmt]
\end{minted}

\section{Ambientes}

O ambiente \(\sigma : \mathit{Var} \rightharpoonup \mathit{Val}\)
agora associa variáveis a \textbf{valores} (inteiros ou booleanos).
A notação \(\sigma(x)\), \(\sigma[x \mapsto v]\) e \(\emptyset\) é a
mesma do capítulo anterior.

\section{Semântica Small-Step}

\subsection{Expressões}

Os valores \(v\) são as formas terminais; qualquer expressão que
não seja um valor pode ainda ser reduzida. As regras para adição e
multiplicação são idênticas ao capítulo anterior. As novas regras
tratam das operações booleanas.

\textbf{Variável:}

\[
  \infer[_{\{Var\}}]{(\sigma,\, x) \to (\sigma,\, \sigma(x))}{x \in
  \mathrm{dom}(\sigma)}
\]

\textbf{Redução de subexpressões:}

Para cada operador binário \(\oplus\), as regras (OpL) e (OpR)
reduzem o operando esquerdo primeiro e depois o direito,
preservando o ambiente:

\[
  \begin{array}{cc}
    \infer[_{\{OpL\}}]{(\sigma,\, e_1 \oplus e_2)\to (\sigma,\, e_1'
    \oplus e_2)}{(\sigma, e_1) \to (\sigma, e_1')}
    &
    \infer[_{\{OpR\}}]{(\sigma,\, v_1 \oplus e_2)\to (\sigma,\, v_1
    \oplus e_2')}{(\sigma, e_2) \to (\sigma, e_2')}
  \end{array}
\]

\textbf{Operações sobre valores:}

\[
  \begin{array}{cc}
    \infer[_{\{Add\}}]{(\sigma,\, n_1 + n_2) \to (\sigma,\, n_1
        \mathbin{\hat{+}}
    n_2)}{}
    &
    \infer[_{\{Mul\}}]{(\sigma,\, n_1 \star n_2) \to (\sigma,\, n_1
    \mathbin{\hat{\star}} n_2)}{}
  \end{array}
\]

\[
  \begin{array}{cc}
    \infer[_{\{Lt\}}]{(\sigma,\, n_1 < n_2) \to (\sigma,\, n_1 \mathbin{\hat{<}}
    n_2)}{}
    &
    \infer[_{\{Eq\}}]{(\sigma,\, v_1 = v_2) \to (\sigma,\, v_1 \mathbin{\hat{=}}
    v_2)}{}
  \end{array}
\]

\[
  \begin{array}{cc}
    \infer[_{\{Or\}}]{(\sigma,\, b_1 \mathbin{|} b_2) \to (\sigma,\, b_1
    \mathbin{\hat{\lor}} b_2)}{}
    &
    \infer[_{\{Not\}}]{(\sigma,\, \neg b) \to (\sigma,\, \hat{\neg}\, b)}{}\\ \\
    \multicolumn{2}{c}{
      \infer[_{\{NotStep\}}]{(\sigma,\, \neg e) \to (\sigma,\, \neg
      e')}{(\sigma,\, e) \to (\sigma,\, e')}
    }
  \end{array}
\]

Os operadores com acento (\(\hat{+}\), \(\hat{<}\), \(\hat{\lor}\),
etc.) denotam as operações matemáticas correspondentes, para
distingui-las dos operadores da linguagem.

\subsection{Comandos}

As regras para \haskell{SAssign},
\haskell{SRead}, \haskell{SPrint} e
para sequências são as mesmas do capítulo anterior. As novas regras
tratam de \haskell{if} e \haskell{while}.

\textbf{Condicional: avaliação da condição}

\[
  \infer[_{\{IfStep\}}]{(\sigma,\; \mathbf{if}\; e\;
    \mathbf{then}\; B_1\; \mathbf{else}\; B_2) \to (\sigma,\;
      \mathbf{if}\; e'\; \mathbf{then}\; B_1\; \mathbf{else}\;
  B_2)}{(\sigma, e) \to (\sigma, e')}
\]

\textbf{Condicional: ramo verdadeiro}

\[
  \infer[_{\{IfTrue\}}]{(\sigma,\; \mathbf{if}\; \mathbf{true}\; \mathbf{then}\;
  B_1\; \mathbf{else}\; B_2) \to (\sigma,\; B_1)}{}
\]

\textbf{Condicional: ramo falso}

\[
  \infer[_{\{IfFalse\}}]{(\sigma,\; \mathbf{if}\; \mathbf{false}\;
      \mathbf{then}\;
  B_1\; \mathbf{else}\; B_2) \to (\sigma,\; B_2)}{}
\]

\textbf{While: iteração}

A regra chave para while é a de
\textbf{iteração condicional}: o laço se transforma em um condicional que,
no ramo verdadeiro, executa o corpo e depois o laço novamente.

\[
  \infer[_{\{While\}}]{(\sigma,\; \mathbf{while}\; e\; \mathbf{do}\; B) \to
    (\sigma,\; \mathbf{if}\; e\; \mathbf{then}\; (B;\; \mathbf{while}\;
  e\; \mathbf{do}\; B)\; \mathbf{else}\; \epsilon)}{}
\]

onde \(\epsilon\) denota a sequência vazia. Esta regra captura a
semântica de repetição de forma simples: o
while se reduz a um condicional cujo ramo
verdadeiro contém o corpo seguido do próprio while.

\subsection{Exemplo: Small-Step}

Considere o programa:

\begin{verbatim}
i := 0
while i < 3 do { i := i + 1 }
print i
\end{verbatim}

Partindo de \(\emptyset\), após a atribuição inicial temos \(\sigma_0
= [i \mapsto 0]\). Mostramos os passos de uma iteração do laço:

\[
  (\sigma_0,\; \mathbf{while}\; i < 3\; \mathbf{do}\; \{i
  \mathbin{:=} i+1\})
\] \[
  \xrightarrow{\text{While}}
  (\sigma_0,\; \mathbf{if}\; i < 3\; \mathbf{then}\; \{i \mathbin{:=}
  i+1;\; \mathbf{while}\ldots\}\; \mathbf{else}\; \epsilon)
\] \[
  \xrightarrow{\text{IfStep, Var, Lt}}
  (\sigma_0,\; \mathbf{if}\; \mathbf{true}\; \mathbf{then}\; \{i
  \mathbin{:=} i+1;\; \mathbf{while}\ldots\}\; \mathbf{else}\; \epsilon)
\] \[
  \xrightarrow{\text{IfTrue}}
  (\sigma_0,\; i \mathbin{:=} i+1;\; \mathbf{while}\; i < 3\;
  \mathbf{do}\; \{i \mathbin{:=} i+1\})
\] \[
  \xrightarrow{\text{Assign}}
  ([i \mapsto 1],\; \mathbf{while}\; i < 3\; \mathbf{do}\; \{i
  \mathbin{:=} i+1\})
\]

Após três iterações, \(\sigma = [i \mapsto 3]\). Na quarta, a
condição \(3 < 3\) reduz para \(\mathbf{false}\) e a regra (IfFalse)
seleciona \(\epsilon\), terminando o laço.

\subsection{Extensão TImp: Semântica Small-Step}\label{sec:timp-small-step}

A linguagem \emph{TImp} estende \emph{While} com construções de
registro e funções. As regras small-step para registros seguem o
mesmo padrão das construções já apresentadas; a semântica completa de
chamadas de função é especificada na forma big-step na
Seção~\ref{sec:timp-big-step}, pois o comando \haskell{return}
exigiria um modelo de pilha de chamadas para ser tratado em small-step.

As configurações passam a incluir o ambiente de funções \(\Phi\):
\((\sigma, \Phi, e)\) para expressões e \((\sigma, \Phi, s)\) para
comandos.

\paragraph{Construção de registro.}

Os campos são avaliados da esquerda para a direita. Enquanto houver
um campo não-valor, ele é reduzido (NewStep); quando todos os campos
são valores, o registro é construído (New):

\[
  \infer[_{\{NewStep\}}]
  {(\sigma, \Phi,\;
    \mathbf{new}\; R\,\{f_1\!=\!v_1,\ldots,f_i\!=\!e_i,\ldots\})
    \to (\sigma', \Phi,\;
  \mathbf{new}\; R\,\{f_1\!=\!v_1,\ldots,f_i\!=\!e_i',\ldots\})}
  {(\sigma, \Phi,\, e_i) \to (\sigma', \Phi,\, e_i')}
\]

\[
  \infer[_{\{New\}}]
  {(\sigma, \Phi,\; \mathbf{new}\; R\,\{f_1\!=\!v_1,\ldots,f_n\!=\!v_n\})
    \to (\sigma, \Phi,\;
  \mathbf{rec}(R,\,\{f_1 \mapsto v_1,\ldots,f_n \mapsto v_n\}))}
  {}
\]

\paragraph{Acesso a campo.}

\[
  \infer[_{\{FieldStep\}}]
  {(\sigma, \Phi,\; e.f) \to (\sigma', \Phi,\; e'.f)}
  {(\sigma, \Phi,\, e) \to (\sigma', \Phi,\, e')}
\]

\[
  \infer[_{\{Field\}}]
  {(\sigma, \Phi,\; \mathbf{rec}(R,\mathit{flds}).f) \to
  (\sigma, \Phi,\; \mathit{flds}(f))}
  {f \in \mathrm{dom}(\mathit{flds})}
\]

\paragraph{Declaração de variável.}

\[
  \infer[_{\{DeclStep\}}]
  {(\sigma, \Phi,\; \mathbf{var}\; x \mathbin{:} \tau = e) \to
  (\sigma', \Phi,\; \mathbf{var}\; x \mathbin{:} \tau = e')}
  {(\sigma, \Phi,\, e) \to (\sigma', \Phi,\, e')}
\]

\[
  \infer[_{\{Decl\}}]
  {(\sigma, \Phi,\; \mathbf{var}\; x \mathbin{:} \tau = v) \to
  (\sigma[x \mapsto v],\; \Phi,\; \epsilon)}
  {}
\]

\paragraph{Atribuição de campo.}

\[
  \infer[_{\{FieldAssignStep\}}]
  {(\sigma, \Phi,\; x.f \mathbin{:=} e) \to
  (\sigma', \Phi,\; x.f \mathbin{:=} e')}
  {(\sigma, \Phi,\, e) \to (\sigma', \Phi,\, e')}
\]

\[
  \infer[_{\{FieldAssign\}}]
  {(\sigma, \Phi,\; x.f \mathbin{:=} v)
    \to (\sigma[x \mapsto \mathbf{rec}(R,\,\mathit{flds}[f \mapsto
  v])],\; \Phi,\; \epsilon)}
  {\sigma(x) = \mathbf{rec}(R,\mathit{flds})}
\]

\paragraph{Avaliação de argumentos.}

Os argumentos de uma chamada de função são avaliados da esquerda para
a direita antes da invocação. A regra de invocação — que cria um
ambiente de ativação fresco e executa o corpo — é especificada na
forma big-step:

\[
  \infer[_{\{CallStep\}}]
  {(\sigma, \Phi,\;
    f(v_1,\ldots,v_{i-1},e_i,\ldots,e_n)) \to
    (\sigma', \Phi,\;
  f(v_1,\ldots,v_{i-1},e_i',\ldots,e_n))}
  {(\sigma, \Phi,\, e_i) \to (\sigma', \Phi,\, e_i')}
\]

\section{Semântica Big-Step}

\subsection{Expressões}

A relação \(\sigma \vdash e \Downarrow v\) agora produz um
\textbf{valor} (inteiro ou booleano).

\[
  \begin{array}{cc}
    \infer[_{\{Val\}}]{\sigma \vdash v \Downarrow v}{}
    &
    \infer[_{\{Var\}}]{\sigma \vdash x
    \Downarrow \sigma(x)}{x \in \mathrm{dom}(\sigma)}
  \end{array}
\]

\[
  \begin{array}{cc}
    \infer[_{\{Add\}}]{\sigma \vdash e_1 + e_2 \Downarrow n_1
    \mathbin{\hat{+}} n_2}{\sigma \vdash e_1 \Downarrow n_1 & \sigma \vdash e_2
    \Downarrow n_2}
    &
    \infer[_{\{Mul\}}]{\sigma \vdash e_1 \star e_2 \Downarrow n_1
    \mathbin{\hat{\star}} n_2}{\sigma \vdash e_1 \Downarrow n_1 &
    \sigma \vdash e_2\Downarrow n_2}
  \end{array}
\]

\[
  \begin{array}{cc}
    \infer[_{\{Lt\}}]{\sigma \vdash e_1 < e_2 \Downarrow n_1
    \mathbin{\hat{<}} n_2}
    {\sigma \vdash e_1 \Downarrow n_1 \quad \sigma \vdash e_2
    \Downarrow n_2}
    &
    \infer[_{\{Eq\}}]{\sigma \vdash e_1 = e_2 \Downarrow v_1
    \mathbin{\hat{=}} v_2}{\sigma \vdash e_1 \Downarrow v_1 & \sigma \vdash e_2
    \Downarrow v_2}
  \end{array}
\]

\[
  \begin{array}{cc}
    \infer[_{\{Or\}}]{\sigma \vdash e_1 \mathbin{|} e_2 \Downarrow b_1
    \mathbin{\hat{\lor}} b_2}{\sigma \vdash e_1 \Downarrow b_1 &
      \sigma \vdash e_2
    \Downarrow b_2}
    &
    \infer[_{\{Not\}}]{\sigma \vdash \neg e \Downarrow \hat{\neg}\, b}
    {\sigma \vdash e \Downarrow b}
  \end{array}
\]

\subsection{Comandos}

As regras para \haskell{SAssign},
\haskell{SRead}, \haskell{SPrint},
(SeqNil) e (SeqCons) permanecem iguais ao capítulo anterior. As novas
regras são:

\textbf{Condicional: ramo verdadeiro}

\[
  \infer[_{\{IfTrue\}}]{\sigma,\; \mathbf{if}\; e\; \mathbf{then}\;
    B_1\; \mathbf{else}\;
  B_2 \;\Downarrow\; \sigma'}{\sigma \vdash e \Downarrow
    \mathbf{true} \quad \sigma,\; B_1
  \Downarrow \sigma'}
\]

\textbf{Condicional: ramo falso}

\[
  \infer[_{\{IfFalse\}}]{\sigma,\; \mathbf{if}\; e\; \mathbf{then}\;
    B_1\; \mathbf{else}\;
  B_2 \;\Downarrow\; \sigma'}{\sigma \vdash e \Downarrow
    \mathbf{false} \quad \sigma,\; B_2
  \Downarrow \sigma'}
\]

\textbf{While: condição falsa (terminação)}

\[
  \infer[_{\{WhileFalse\}}]{\sigma,\; \mathbf{while}\; e\;
  \mathbf{do}\; B \;\Downarrow\; \sigma}
  {\sigma \vdash e \Downarrow \mathbf{false}}
\]

\textbf{While: condição verdadeira (iteração)}

\[
  \infer[_{\{WhileTrue\}}]{\sigma,\; \mathbf{while}\; e\;
  \mathbf{do}\; B \;\Downarrow\; \sigma''}
  {\sigma \vdash e \Downarrow \mathbf{true}
    &  \sigma,\; B \Downarrow \sigma'
  & \sigma',\; \mathbf{while}\; e\; \mathbf{do}\; B \Downarrow \sigma''}
\]

A regra (WhileTrue) é recursiva: após executar o corpo \(B\) e obter
\(\sigma'\), o while é aplicado novamente a
partir de \(\sigma'\). A recursão termina quando a condição se torna
\(\mathbf{false}\), coberta por (WhileFalse). Se a condição nunca se
tornar falsa, não existe derivação big-step finita: o programa diverge.

\subsection{Exemplos: Big-Step}

\textbf{Exemplo 1:} Considere o seguinte programa:

\begin{verbatim}
if 2 < 3 then {
  x := 10;
} else {
  x := 20;
}
\end{verbatim}

\[
  \dfrac{
    \dfrac{
      \dfrac{}{\emptyset \vdash 2 \Downarrow 2}
      \quad
      \dfrac{}{\emptyset \vdash 3 \Downarrow 3}
    }{\emptyset \vdash 2 < 3 \Downarrow \mathbf{true}}
    \quad
    \dfrac{
      \dfrac{}{\emptyset \vdash 10 \Downarrow 10}
    }{\emptyset,\; x \mathbin{:=} 10 \;\Downarrow\; [x \mapsto 10]}
  }{\emptyset,\; \mathbf{if}\; 2<3\; \mathbf{then}\; \{x \mathbin{:=}
    10\}\; \mathbf{else}\; \{x \mathbin{:=} 20\} \;\Downarrow\; [x
  \mapsto 10]}
\]

\textbf{Exemplo 2:} Considere o programa:

\begin{verbatim}
while i < 2 do {
   i := i + 1;
}
\end{verbatim}
e o estado inicial \(\sigma_0 = [i \mapsto 0]\).

Abreviemos \(\sigma_k = [i \mapsto k]\) e o programa anterior
por \(W\).

\[
  \dfrac{
    \sigma_0 \vdash i < 2 \Downarrow \mathbf{true}
    \quad
    \sigma_0,\; i \mathbin{:=} i+1 \;\Downarrow\; \sigma_1
    \quad
    \dfrac{
      \sigma_1 \vdash i < 2 \Downarrow \mathbf{true}
      \quad
      \sigma_1,\; i \mathbin{:=} i+1 \;\Downarrow\; \sigma_2
      \quad
      \dfrac{
        \sigma_2 \vdash i < 2 \Downarrow \mathbf{false}
      }{\sigma_2,\; W \;\Downarrow\; \sigma_2}
    }{\sigma_1,\; W \;\Downarrow\; \sigma_2}
  }{\sigma_0,\; W \;\Downarrow\; \sigma_2}
\]

A derivação é uma árvore de profundidade proporcional ao número de
iterações. O laço termina porque \(i\) cresce a cada iteração e
eventualmente atinge 2.

\subsection{Extensão TImp: Semântica Big-Step}\label{sec:timp-big-step}

A linguagem \emph{TImp} estende \emph{While} com declarações de
funções e tipos registro. Nesta seção apresentamos os domínios
semânticos e as regras big-step para as novas construções.

\subsubsection{Sintaxe Estendida}

Um programa TImp consiste em uma sequência de declarações seguida de
um bloco de comandos de nível superior:

\[
  \begin{array}{lcl}
    \mathit{prog}  & ::= & \mathit{decl}^*\; s^* \\[4pt]
    \mathit{decl}  & ::= &
    \mathbf{record}\; R\; \{\, \overline{f_i \mathbin{:} \tau_i}\,\}  \\
    & \mid &
    \mathbf{fn}\; f\,(\overline{x_i \mathbin{:} \tau_i}) \mathbin{:}
    \rho\; \{\, s^*\,\} \\[4pt]
    \rho           & ::= & \mathbf{void} \mid \tau \\[4pt]
    \tau           & ::= & \mathbf{int} \mid \mathbf{bool} \mid
    \mathbf{string} \mid R \\[4pt]
    s              & ::= & \cdots \mid
    \mathbf{var}\; x \mathbin{:} \tau \mathbin{=} e \mid
    x.f \mathbin{:=} e \mid
    \mathbf{return}\; e \mid \mathbf{return} \mid
    f(e_1,\ldots,e_n) \\[4pt]
    e              & ::= & \cdots \mid e.f \mid
    \mathbf{new}\; R\,\{\,\overline{f_i \mathbin{=} e_i}\,\} \mid
    f(e_1,\ldots,e_n)
  \end{array}
\]

Os reticências \(\cdots\) denotam as construções já presentes em
\emph{While}.

\subsubsection{Domínios Semânticos Estendidos}

\paragraph{Valores.}
O domínio de valores é estendido com valores registro e com o valor
unitário (resultado de funções \haskell{void}):

\[
  v \;::=\; n \;\mid\; b \;\mid\; s \;\mid\;
  \mathbf{rec}(R,\, \{f_1 \mapsto v_1, \ldots, f_n \mapsto v_n\})
  \;\mid\; \mathbf{unit}
\]

Um valor registro \(\mathbf{rec}(R, \mathit{flds})\) carrega o nome do
tipo \(R\) e uma função parcial \(\mathit{flds}\) que mapeia nomes de
campo a valores. A semântica adota \textbf{semântica de valor} para
registros: toda atribuição e passagem de parâmetro copia o registro
inteiro.

\paragraph{Ambientes de funções.}
Além do ambiente de variáveis \(\sigma\), a semântica introduz um
\textbf{ambiente de funções} global

\[
  \Phi : \mathit{Name} \rightharpoonup
  (\overline{x_i},\; B)
\]

que mapeia cada nome de função à sua lista de parâmetros formais
\(\overline{x_i}\) e ao seu corpo \(B\). O ambiente \(\Phi\) é
construído uma única vez a partir das declarações do programa e
permanece imutável durante a execução.

\paragraph{Resultado de bloco.}
As construções \haskell{return} interrompem a execução normal de um
bloco e devolvem um valor ao chamador. Para modelar esse comportamento,
a relação de avaliação de comandos e blocos passa a produzir um
\textbf{resultado}:

\[
  r \;::=\; \mathbf{normal} \;\mid\; {\uparrow}\, v
\]

O resultado \(\mathbf{normal}\) indica terminação sem retorno; o
resultado \({\uparrow}\, v\) indica que um \haskell{return} foi
executado e o valor \(v\) está sendo propagado para o ponto de chamada.
A relação de avaliação de comandos torna-se \(\sigma,\Phi \vdash s
\Downarrow (\sigma', r)\) e a de blocos, \(\sigma,\Phi \vdash B
\Downarrow (\sigma', r)\).

\subsubsection{Semântica Big-Step: Registros}

\paragraph{Expressões}

\textbf{Construção de registro:}

\[
  \infer[_{\{\mathit{New}\}}]
  {\sigma,\Phi \vdash \mathbf{new}\; R\,\{f_1\!=\!e_1,\ldots,f_n\!=\!e_n\}
  \Downarrow \mathbf{rec}(R,\,\{f_1 \mapsto v_1,\ldots,f_n \mapsto v_n\})}
  {\sigma,\Phi \vdash e_1 \Downarrow v_1
    \quad \cdots \quad
  \sigma,\Phi \vdash e_n \Downarrow v_n}
\]

Cada campo é avaliado na ordem em que aparece na expressão.

\textbf{Acesso a campo:}

\[
  \infer[_{\{\mathit{Field}\}}]
  {\sigma,\Phi \vdash e.f \Downarrow \mathit{flds}(f)}
  {\sigma,\Phi \vdash e \Downarrow \mathbf{rec}(R,\mathit{flds})
  \quad f \in \mathrm{dom}(\mathit{flds})}
\]

O acesso avalia \(e\) para um valor registro e extrai o campo \(f\).
Se \(f\) não existir, a execução trava (erro de tipo prevenido pela
verificação de tipos).

\paragraph{Comandos}

\textbf{Declaração de variável:}

\[
  \infer[_{\{\mathit{Decl}\}}]
  {\sigma,\Phi \vdash \mathbf{var}\; x \mathbin{:} \tau \mathbin{=} e
  \Downarrow (\sigma[x \mapsto v],\; \mathbf{normal})}
  {\sigma,\Phi \vdash e \Downarrow v}
\]

\textbf{Atribuição de campo:}

\[
  \infer[_{\{\mathit{FieldAssign}\}}]
  {\sigma,\Phi \vdash x.f \mathbin{:=} e
    \Downarrow (\sigma[x \mapsto \mathbf{rec}(R,\,\mathit{flds}[f
      \mapsto v])],\;
  \mathbf{normal})}
  {\sigma(x) = \mathbf{rec}(R,\mathit{flds})
  \quad \sigma,\Phi \vdash e \Downarrow v}
\]

A atribuição não aloca um novo registro: ela atualiza o registro em
\(x\) substituindo o campo \(f\) pelo novo valor \(v\). Como a
semântica é de valor, \(x\) recebe uma cópia modificada.

\subsubsection{Semântica Big-Step: Funções}

\paragraph{Propagação de \texttt{return}}

Para que \haskell{return} interrompa a execução de um bloco, as
regras de sequência precisam propagar o resultado.

\textbf{Bloco vazio:}

\[
  \infer[_{\{\mathit{BlkNil}\}}]
  {\sigma,\Phi \vdash \epsilon \Downarrow (\sigma,\; \mathbf{normal})}
  {}
\]

\textbf{Bloco com retorno antecipado:}

\[
  \infer[_{\{\mathit{BlkRet}\}}]
  {\sigma,\Phi \vdash s\,;\,B \Downarrow (\sigma',\; {\uparrow}\,v)}
  {\sigma,\Phi \vdash s \Downarrow (\sigma',\; {\uparrow}\,v)}
\]

Se o primeiro comando produz um retorno, o resto do bloco \(B\) não é
executado e o resultado é propagado imediatamente.

\textbf{Bloco com continuação normal:}

\[
  \infer[_{\{\mathit{BlkCons}\}}]
  {\sigma,\Phi \vdash s\,;\,B \Downarrow r}
  {\sigma,\Phi \vdash s \Downarrow (\sigma',\; \mathbf{normal})
  \quad \sigma',\Phi \vdash B \Downarrow r}
\]

\paragraph{Comandos de retorno}

\textbf{Retorno com valor:}

\[
  \infer[_{\{\mathit{RetVal}\}}]
  {\sigma,\Phi \vdash \mathbf{return}\; e \Downarrow (\sigma,\; {\uparrow}\,v)}
  {\sigma,\Phi \vdash e \Downarrow v}
\]

\textbf{Retorno sem valor (\texttt{void}):}

\[
  \infer[_{\{\mathit{RetVoid}\}}]
  {\sigma,\Phi \vdash \mathbf{return} \Downarrow (\sigma,\;
  {\uparrow}\,\mathbf{unit})}
  {}
\]

\paragraph{Chamada de função}

\textbf{Chamada como expressão} (função não-\texttt{void}):

\[
  \infer[_{\{\mathit{CallExpr}\}}]
  {\sigma,\Phi \vdash f(e_1,\ldots,e_n) \Downarrow v}
  {
    \Phi(f) = (\overline{x_i},\; B)
    \quad
    \sigma,\Phi \vdash e_i \Downarrow v_i \;\;(\forall\, i)
    \quad
    [x_1 \mapsto v_1,\ldots,x_n \mapsto v_n],\Phi \vdash B
    \Downarrow (\_,\; {\uparrow}\,v)
  }
\]

O corpo de \(f\) é executado em um \textbf{ambiente fresco}
\([x_1 \mapsto v_1,\ldots,x_n \mapsto v_n]\): apenas os parâmetros
formais estão em escopo. O resultado \({\uparrow}\,v\) fornece o valor
de retorno; o ambiente da chamada \(\sigma\) não é modificado.

\textbf{Chamada como comando} (função \texttt{void} ou retorno descartado):

\[
  \infer[_{\{\mathit{CallStmt}\}}]
  {\sigma,\Phi \vdash f(e_1,\ldots,e_n) \Downarrow (\sigma,\; \mathbf{normal})}
  {
    \Phi(f) = (\overline{x_i},\; B)
    \quad
    \sigma,\Phi \vdash e_i \Downarrow v_i \;\;(\forall\, i)
    \quad
    [x_1 \mapsto v_1,\ldots,x_n \mapsto v_n],\Phi \vdash B \Downarrow (\_,\; \_)
  }
\]

O ambiente de chamada é devolvido inalterado: a função não pode
alterar o estado externo através de variáveis globais mutáveis nem por
passagem por referência.

\section{Implementação em Haskell}

\subsection{Expressões}

A função \haskell{interpExp} retorna agora um
\haskell{Value} (inteiro ou booleano). As novas
equações implementam diretamente as regras semânticas:

\begin{minted}{haskell}
interpExp :: Exp -> Interp Value
interpExp (EValue v) = pure v
interpExp (EVar v) = askVar v
interpExp (e1 :+: e2) = do { v1 <- interpExp e1; v2 <- interpExp e2; v1 .+. v2 }
interpExp (e1 :*: e2) = do { v1 <- interpExp e1; v2 <- interpExp e2; v1 .*. v2 }
interpExp (e1 :<: e2) = do { v1 <- interpExp e1; v2 <- interpExp e2; v1 .<. v2 }
interpExp (e1 :=: e2) = do { v1 <- interpExp e1; v2 <- interpExp e2; v1 .=. v2 }
interpExp (e1 :|: e2) = do { v1 <- interpExp e1; v2 <- interpExp e2; v1 .|. v2 }
interpExp (ENot e) = do { v <- interpExp e; vnot v }
\end{minted}

As funções \haskell{(.+.)},
\haskell{(.<.)}, \haskell{(.=.)},
\haskell{(.|.)} e \haskell{vnot},
definidas em \haskell{Utils.Value}, verificam que os
operandos têm os tipos corretos e lançam um erro via
\haskell{throwError} caso contrário.

\subsection{Condicional}

A equação para \haskell{SIf} implementa as regras
(IfTrue) e (IfFalse):

\begin{minted}{haskell}
interpStmt (SIf e blk1 blk2) = do
    v <- interpExp e
    case v of
      VBool True  -> mapM_ interpStmt blk1
      VBool False -> mapM_ interpStmt blk2
      _           -> throwError $ unlines
                       ["Expecting a boolean value, but found:", pretty e]
\end{minted}

O casamento de padrões em \haskell{VBool True} e
\haskell{VBool False} corresponde diretamente às
premissas das regras (IfTrue) e (IfFalse). O caso
\haskell{_} trata erros de tipo que as regras
semânticas simplesmente não cobrem.

\subsection{Repetição}

A equação para \haskell{SWhile} implementa as regras
(WhileFalse) e (WhileTrue):

\begin{minted}{haskell}
interpStmt (SWhile e blk1) = do
    v <- interpExp e
    case v of
      VBool False -> pure ()
      VBool True  -> do
          mapM_ interpStmt blk1
          interpStmt (SWhile e blk1)
      _ -> throwError $ unlines
             ["Expecting a boolean value, but found:", pretty e]
\end{minted}

A correspondência com as regras semânticas é precisa:

\begin{itemize}
  \item
    \haskell{VBool False -> pure ()}, correspondendo à regra
    (WhileFalse): o ambiente não muda, a execução termina sem
    produzir nenhum efeito adicional.
  \item
    \haskell{VBool True -> mapM_ interpStmt blk1 >>
    interpStmt (SWhile e blk1)}, correspondendo à regra (WhileTrue): executa o
    bloco (\(\sigma, B \Downarrow \sigma'\)) e depois invoca
    \haskell{interpStmt (SWhile e blk1)}
    recursivamente (\(\sigma', \mathbf{while}\ldots \Downarrow \sigma''\)).
\end{itemize}

A recursão na implementação espelha a recursão na regra semântica: se
o laço não terminar, a função Haskell também não terminará.

\subsection{Executando o Interpretador}

\begin{verbatim}
cabal run while -- --interp --file programa.while
\end{verbatim}

Por exemplo, o seguinte programa calcula o fatorial de 5:

\begin{verbatim}
n := 5;
acc := 1;
while 0 < n do {
    acc := acc * n;
    n := n + (-1);
}
print acc;
\end{verbatim}

O interpretador imprime 120.

\subsection{Implementação da Extensão TImp}

A implementação está em \verb|TImp/Interp/TImpInterp.hs|. As regras
semânticas apresentadas na Seção~\ref{sec:timp-big-step} têm
correspondência direta com o código Haskell.

\paragraph{Domínio de valores e estado.}

\begin{minted}{haskell}
data Value
  = VInt    Int
  | VBool   Bool
  | VString String
  | VRecord Name (Map Field Value)  -- rec(R, flds)
  | VUnit                           -- unit

data InterpState = InterpState
  { varEnv  :: Map Var Value        -- sigma
  , funcEnv :: Map Name FuncDecl    -- Phi
  }

type InterpM a = StateT InterpState (ExceptT String IO) a
\end{minted}

O campo \haskell{funcEnv} corresponde ao ambiente global \(\Phi\),
construído uma vez no ponto de entrada e nunca modificado.

\paragraph{Resultado de bloco.}

A função \haskell{execBlock} retorna \haskell{Maybe Value},
onde \haskell{Nothing} corresponde a \(\mathbf{normal}\) e
\haskell{Just v} corresponde a \({\uparrow}\,v\):

\begin{minted}{haskell}
execBlock :: [Stmt] -> InterpM (Maybe Value)
execBlock []     = pure Nothing              -- BlkNil
execBlock (s:ss) = do
  r <- execStmt s
  case r of
    Just _  -> pure r                        -- BlkRet: propagate
    Nothing -> execBlock ss                  -- BlkCons: continue
\end{minted}

\paragraph{Registros.}

A avaliação de \haskell{ENew} implementa a regra (New):

\begin{minted}{haskell}
evalExp (ENew rname initFields) = do
  vals <- mapM (\(f, e) -> (f,) <$> evalExp e) initFields
  pure (VRecord rname (Map.fromList vals))
\end{minted}

O acesso a campo implementa a regra (Field):

\begin{minted}{haskell}
evalExp (EField e f) = do
  v <- evalExp e
  case v of
    VRecord _ fields ->
      case Map.lookup f fields of
        Just val -> pure val
        Nothing  -> throwError ("Record has no field: " ++ f)
    _ -> throwError "Field access on non-record value"
\end{minted}

A atribuição de campo implementa (FieldAssign):

\begin{minted}{haskell}
execStmt (SFieldAssign v f e) = do
  rval <- lookupVar v
  newVal <- evalExp e
  case rval of
    VRecord name fields ->
      let fields' = Map.insert f newVal fields
      in  modify (\st -> st
            { varEnv = Map.insert v (VRecord name fields') (varEnv st) })
    _ -> throwError (v ++ " is not a record")
  pure Nothing
\end{minted}

\paragraph{Funções.}

O comando \haskell{return} implementa as regras (RetVal) e (RetVoid):

\begin{minted}{haskell}
execStmt (SReturn Nothing)  = pure (Just VUnit)
execStmt (SReturn (Just e)) = Just <$> evalExp e
\end{minted}

A chamada de função implementa as regras (CallExpr) e (CallStmt).
A função auxiliar \haskell{callFunc} cria um novo ambiente e executa
o corpo:

\begin{minted}{haskell}
callFunc :: Name -> [Value] -> InterpM Value
callFunc fname argVals = do
  fenv <- gets funcEnv
  case Map.lookup fname fenv of
    Nothing -> throwError ("Undefined function: " ++ fname)
    Just (FuncDecl _ params _ body) -> do
      let bindings = zip [v | Param v _ <- params] argVals
      withFreshEnv bindings $ do        -- ambiente novo: sigma_0
        r <- execBlock body
        pure (fromMaybe VUnit r)        -- extrai v de (uparrow v)
\end{minted}

A função \haskell{withFreshEnv} salva o ambiente corrente, inicializa
o novo ambiente com apenas os parâmetros formais e restaura o ambiente
original ao terminar, implementando a semântica de escopo local:

\begin{minted}{haskell}
withFreshEnv :: [(Var, Value)] -> InterpM a -> InterpM a
withFreshEnv bindings m = do
  saved <- gets varEnv
  modify (\st -> st { varEnv = Map.fromList bindings })
  r <- m
  modify (\st -> st { varEnv = saved })
  pure r
\end{minted}

\section{Conclusão}

Neste capítulo, completámos a semântica operacional de uma linguagem
imperativa com as construções de controle de fluxo mais fundamentais:

\begin{itemize}
  \item
    O \textbf{condicional} é modelado por duas regras simétricas,
    (IfTrue) e (IfFalse), que selecionam o ramo a executar com base
    no valor booleano da condição.
  \item
    O \textbf{laço} while é modelado por
    (WhileFalse), que termina quando a condição é falsa, e
    (WhileTrue), que executa o corpo e recomeça: uma regra
    inerentemente recursiva que espelha a definição operacional
    de repetição.
  \item
    Na semântica small-step, o while se
    \textbf{desenrola} em um condicional, o que torna explícita a
    estrutura de cada iteração.
  \item
    Na semântica big-step, a não-terminação corresponde à ausência de
    uma derivação finita: a semântica é parcial por construção.
  \item
    \textbf{Registros} são modelados como valores estruturados
    \(\mathbf{rec}(R, \mathit{flds})\) com semântica de cópia; as regras
    (New), (Field) e (FieldAssign) cobrem criação, leitura e atualização.
  \item
    \textbf{Funções} introduzem o conceito de \textit{resultado de bloco}
    \(r \in \{\mathbf{normal}, {\uparrow}\,v\}\) para modelar
    \haskell{return}: as regras (BlkRet) e (BlkCons) propagam ou
    descartam esse resultado; (CallExpr) executa a função em um ambiente
    fresco e extrai o valor retornado.
\end{itemize}

A implementação em \verb|TImp/Interp/TImpInterp.hs| segue a
especificação semântica com precisão, confirmando que uma boa
especificação formal é também um guia de implementação direta.
